%% kapa: exam=1 round=20 year=2009 period=2 form=A subject=appraisal_law subject_ko=부동산관계법규 seq=2 number=42
%% source: Q-Net 감정평가사 기출문제 (cst00309) artlSeq=5024178 file=extracted/1008/5024178/2009년 제20회 감정평가사 1차시험 기출문제/2교시.hwp
%% conversion: kordoc 4.24.51 -> tools/gen1.py (auto); review=proofread-sonnet; flags=none
\begin{question}{42}
국토의 계획 및 이용에 관한 법령상 용도지역ㆍ용도지구 및 용도구역에 관한 설명으로 옳지 \uline{않은} 것은?
\begin{choices}
  \item 용도지역은 도시지역ㆍ관리지역ㆍ농림지역ㆍ자연환경보전지역으로 구분되며, 도시지역은 주거지역ㆍ상업지역ㆍ공업지역ㆍ녹지지역으로 구분된다.
  \item 시ㆍ도의 도시계획조례로 국토의 계획 및 이용에 관한 법률상 용도지구 외의 용도지구를 정할 때에는 당해 용도지역 또는 용도구역의 행위제한을 완화 또는 강화하는 용도지구를 신설해서는 아니 된다.
  \item 국토해양부장관은 5년 이상 20년 이내의 기간 동안 시가화를 유보할 필요가 있다고 인정되는 경우 시가화조정구역의 지정을 도시관리계획으로 결정할 수 있다.
  \item 도시지역으로의 편입이 예상되는 지역으로서 계획적ㆍ체계적인 관리가 필요한 지역은 계획관리지역에 해당한다.
  \item 집단취락지구는 녹지지역ㆍ관리지역ㆍ농림지역 또는 자연환경보전지역 안의 취락을 정비하기 위하여 필요한 지구를 말한다.
\end{choices}
\end{question}
\answer{2,5}
\answernote{복수 정답(정답표 원문: 2, 5)}
